% =========================================================
% Section: IEEE 754 Floating-Point Formats
% =========================================================

\subsection*{IEEE 754 Floating-Point Formats}

\providecommand{\IEEEFloatLayout}[8]{%
\begin{center}
\begin{tikzpicture}[x=#1 cm,y=0.8cm]
    \pgfmathsetmacro{\SignEnd}{#2}
    \pgfmathsetmacro{\ExpEnd}{#2 + #3}
    \pgfmathsetmacro{\FracEnd}{#2 + #3 + #4}
    \pgfmathtruncatemacro{\SignBitLabel}{#5 - 1}

    \draw (0,0) rectangle (\SignEnd,1);
    \draw (\SignEnd,0) rectangle (\ExpEnd,1);
    \draw (\ExpEnd,0) rectangle (\FracEnd,1);

    \node at ({\SignEnd/2},0.5) {\small $s$};
    \node at ({(\SignEnd+\ExpEnd)/2},0.5) {\small exponent};
    \node at ({(\ExpEnd+\FracEnd)/2},0.5) {\small fraction};

    \node[above] at (0.5,1) {\scriptsize #5};
    \node[above] at ({\SignEnd+0.5},1) {\scriptsize \SignBitLabel};
    \node[above] at ({\ExpEnd-0.5},1) {\scriptsize #6};
    \node[above] at ({\ExpEnd+0.5},1) {\scriptsize #7};
    \node[above] at ({\FracEnd-0.5},1) {\scriptsize 0};
\end{tikzpicture}
\end{center}
}


% =========================================================
% General structure
% =========================================================

\begin{definitionbox}{Definition - IEEE 754 Standard}
\begin{definition}[IEEE 754 Standard]
\label{def:numerical-analysis-ieee-754-standard}
The \emph{IEEE 754 standard} is the standard convention used here for binary
floating-point data, formats, special encodings, and finite-precision numerical
arithmetic.
\end{definition}

\DefinitionalRoot
\end{definitionbox}

\begin{remark*}[Standard quantified statement]
This definition fixes the local standard that governs the IEEE binary formats
and encodings used in this section.
\end{remark*}

\begin{remark*}[Interpretation]
The standard is the organizing object. Concrete formats such as binary32 and
binary64, and concrete encoding classes such as normal, subnormal, infinity,
and NaN, are specializations governed by this standard.
\end{remark*}

\begin{definitionbox}{Definition - IEEE 754 Binary Floating-Point Datum}
\begin{definition}[IEEE 754 Binary Floating-Point Datum]
\label{def:numerical-analysis-ieee-754-binary-floating-point-datum}
An \emph{IEEE 754 binary floating-point datum} is a binary encoding partitioned into
three fields:

\begin{itemize}
    \item a \emph{sign field} consisting of one bit,
    \item an \emph{exponent field} consisting of a fixed number of bits,
    \item a \emph{fraction field} consisting of a fixed number of bits.
\end{itemize}

Its interpretation depends on the exponent-field pattern and the fraction-field pattern.
\end{definition}
\end{definitionbox}

\begin{remark*}[Standard quantified statement]
\[
x=\mathsf{IEEEBinaryDatum}(s,e,f,F)
\Longleftrightarrow
\operatorname{SignField}(x)=s
\land
\operatorname{ExponentField}(x)=e
\land
\operatorname{FractionField}(x)=f.
\]
\end{remark*}

\begin{remark*}[Interpretation]
IEEE 754 does not interpret a floating-point bit string as a plain binary integer.
Instead, the sign field determines sign, the exponent field determines scale and
special cases, and the fraction field records the significant bits of the encoded value.
\end{remark*}

\begin{dependencies}
\begin{itemize}
  \item \hyperref[def:numerical-analysis-ieee-754-standard]{IEEE 754 Standard}
\end{itemize}
\end{dependencies}

% =========================================================
% Normalized numbers
% =========================================================

\begin{definitionbox}{Definition - Normalized IEEE 754 Binary Number}
\begin{definition}[Normalized IEEE 754 Binary Number]
\label{def:numerical-analysis-normalized-ieee-754-binary-number}
Let an IEEE 754 binary floating-point format have exponent bias $\mathrm{bias}$.
A finite encoded datum is called \emph{normalized} if its exponent field is neither
all zeros nor all ones.

If $s$ is the sign bit, if $e$ is the unsigned integer determined by the exponent field,
and if $f$ is the binary fraction determined by the fraction field, then the represented
value is
\[
(-1)^s (1.f)_2 2^{e-\mathrm{bias}}.
\]
\end{definition}
\end{definitionbox}

\begin{remark*}[Standard quantified statement]
\[
\operatorname{HasIEEEClass}(x,F,\mathsf{IEEEEncodingClass}(\mathrm{normal}))
\Longleftrightarrow
x=\mathsf{IEEEBinaryDatum}(s,e,f,F)
\land
0<\operatorname{ExponentField}(x)<2^{\operatorname{ExponentWidth}(F)}-1.
\]
\end{remark*}

\begin{remark*}[Interpretation]
In a normalized encoding, the leading binary digit of the significand is implicitly $1$.
This hidden leading bit is what gives normalized numbers one more bit of effective
precision than the stored fraction field alone.
\end{remark*}

\begin{dependencies}
\begin{itemize}
  \item \hyperref[def:numerical-analysis-ieee-754-standard]{IEEE 754 Standard}
  \item \hyperref[def:numerical-analysis-ieee-754-binary-floating-point-datum]{IEEE 754 Binary Floating-Point Datum}
\end{itemize}
\end{dependencies}

% =========================================================
% Subnormal numbers
% =========================================================

\begin{definitionbox}{Definition - Subnormal IEEE 754 Binary Number}
\begin{definition}[Subnormal IEEE 754 Binary Number]
\label{def:numerical-analysis-subnormal-ieee-754-binary-number}
Let an IEEE 754 binary floating-point format have exponent bias $\mathrm{bias}$.
A finite encoded datum is called \emph{subnormal} if its exponent field is all zeros
and its fraction field is not all zeros.

If $s$ is the sign bit and if $f$ is the binary fraction determined by the fraction field,
then the represented value is
\[
(-1)^s (0.f)_2 2^{1-\mathrm{bias}}.
\]
\end{definition}
\end{definitionbox}

\begin{remark*}[Standard quantified statement]
\[
\operatorname{HasIEEEClass}(x,F,\mathsf{IEEEEncodingClass}(\mathrm{subnormal}))
\Longleftrightarrow
x=\mathsf{IEEEBinaryDatum}(s,e,f,F)
\land
\operatorname{ExponentField}(x)=0
\land
\operatorname{FractionField}(x)\ne 0.
\]
\end{remark*}

\begin{remark*}[Interpretation]
Subnormal numbers fill the gap between zero and the smallest positive normalized number.
They preserve gradual underflow by allowing the leading significand bit to be $0$ instead
of the implicit $1$ used for normalized numbers.
\end{remark*}

\begin{dependencies}
\begin{itemize}
  \item \hyperref[def:numerical-analysis-ieee-754-standard]{IEEE 754 Standard}
  \item \hyperref[def:numerical-analysis-ieee-754-binary-floating-point-datum]{IEEE 754 Binary Floating-Point Datum}
\end{itemize}
\end{dependencies}

% =========================================================
% Infinity
% =========================================================

\begin{definitionbox}{Definition - IEEE 754 Infinity Encoding}
\begin{definition}[IEEE 754 Infinity Encoding]
\label{def:numerical-analysis-ieee-754-infinity-encoding}
An IEEE 754 binary floating-point datum encodes \emph{infinity} if its exponent field
is all ones and its fraction field is all zeros.

The sign bit determines whether the represented value is $+\infty$ or $-\infty$.
\end{definition}\end{definitionbox}

\begin{remark*}[Standard quantified statement]
\[
\operatorname{HasIEEEClass}(x,F,\mathsf{IEEEEncodingClass}(\mathrm{infinity}))
\Longleftrightarrow
x=\mathsf{IEEEBinaryDatum}(s,e,f,F)
\land
\operatorname{ExponentField}(x)=2^{\operatorname{ExponentWidth}(F)}-1
\land
\operatorname{FractionField}(x)=0.
\]
\end{remark*}

\begin{remark*}[Interpretation]
Infinity is a distinguished non-finite encoding used to represent overflow and certain
limit-like arithmetic results. The sign field still matters, so IEEE 754 distinguishes
positive infinity from negative infinity.
\end{remark*}

\begin{dependencies}
\begin{itemize}
  \item \hyperref[def:numerical-analysis-ieee-754-standard]{IEEE 754 Standard}
  \item \hyperref[def:numerical-analysis-ieee-754-binary-floating-point-datum]{IEEE 754 Binary Floating-Point Datum}
\end{itemize}
\end{dependencies}

% =========================================================
% NaN
% =========================================================

\begin{definitionbox}{Definition - IEEE 754 NaN Encoding}
\begin{definition}[IEEE 754 NaN Encoding]
\label{def:numerical-analysis-ieee-754-nan-encoding}
An IEEE 754 binary floating-point datum encodes \emph{NaN} (Not a Number) if its
exponent field is all ones and its fraction field is not all zeros.
\end{definition}
\end{definitionbox}

\begin{remark*}[Standard quantified statement]
\[
\operatorname{HasIEEEClass}(x,F,\mathsf{IEEEEncodingClass}(\mathrm{nan}))
\Longleftrightarrow
x=\mathsf{IEEEBinaryDatum}(s,e,f,F)
\land
\operatorname{ExponentField}(x)=2^{\operatorname{ExponentWidth}(F)}-1
\land
\operatorname{FractionField}(x)\ne 0.
\]
\end{remark*}

\begin{remark*}[Interpretation]
A NaN encoding represents an undefined or invalid numerical result rather than a real number.
For example, NaNs arise from operations such as invalid indeterminate forms or from the
propagation of an already invalid datum through later computations.
\end{remark*}

\begin{dependencies}
\begin{itemize}
  \item \hyperref[def:numerical-analysis-ieee-754-standard]{IEEE 754 Standard}
  \item \hyperref[def:numerical-analysis-ieee-754-binary-floating-point-datum]{IEEE 754 Binary Floating-Point Datum}
\end{itemize}
\end{dependencies}

% =========================================================
% Binary16
% =========================================================

\begin{definitionbox}{Definition - IEEE 754 Binary16 Format}
\begin{definition}[IEEE 754 Binary16 Format]
\label{def:numerical-analysis-ieee-754-binary16-format}
The \emph{IEEE 754 binary16 format} is the 16-bit binary floating-point interchange format
with the following field layout:

\begin{itemize}
    \item $1$ sign bit,
    \item $5$ exponent bits,
    \item $10$ fraction bits,
    \item exponent bias $15$.
\end{itemize}

For a normalized finite binary16 datum, the represented value is
\[
(-1)^s (1.f)_2 2^{e-15}.
\]
\end{definition}
\end{definitionbox}

\begin{remark*}[Standard quantified statement]
\[
\mathsf{Binary16Format}
=\mathsf{IEEEBinaryFormat}(16,1,5,10,15).
\]
\end{remark*}

\begin{remark*}[Interpretation]
Binary16 is commonly called \emph{half precision}. It is useful when reduced storage
and bandwidth matter more than wide dynamic range or high precision.
\end{remark*}

\begin{remark*}[Exposition]
\IEEEFloatLayout{0.55}{1}{5}{10}{15}{10}{9}{binary16}
\end{remark*}

\begin{dependencies}
\begin{itemize}
  \item \hyperref[def:numerical-analysis-ieee-754-standard]{IEEE 754 Standard}
  \item \hyperref[def:numerical-analysis-ieee-754-binary-floating-point-datum]{IEEE 754 Binary Floating-Point Datum}
\end{itemize}
\end{dependencies}

% =========================================================
% Binary32
% =========================================================

\begin{definitionbox}{Definition - IEEE 754 Binary32 Format}
\begin{definition}[IEEE 754 Binary32 Format]
\label{def:numerical-analysis-ieee-754-binary32-format}
The \emph{IEEE 754 binary32 format} is the 32-bit binary floating-point interchange format
with the following field layout:

\begin{itemize}
    \item $1$ sign bit,
    \item $8$ exponent bits,
    \item $23$ fraction bits,
    \item exponent bias $127$.
\end{itemize}

For a normalized finite binary32 datum, the represented value is
\[
(-1)^s (1.f)_2 2^{e-127}.
\]
\end{definition}
\end{definitionbox}

\begin{remark*}[Standard quantified statement]
\[
\mathsf{Binary32Format}
=\mathsf{IEEEBinaryFormat}(32,1,8,23,127).
\]
\end{remark*}

\begin{remark*}[Interpretation]
Binary32 is commonly called \emph{single precision}. It is the standard 32-bit floating-point
format used broadly in scientific computing, computer graphics, simulation, and machine learning.
\end{remark*}

\begin{remark*}[Exposition]
\IEEEFloatLayout{0.28}{1}{8}{23}{31}{23}{22}{binary32}
\end{remark*}

\begin{dependencies}
\begin{itemize}
  \item \hyperref[def:numerical-analysis-ieee-754-standard]{IEEE 754 Standard}
  \item \hyperref[def:numerical-analysis-ieee-754-binary-floating-point-datum]{IEEE 754 Binary Floating-Point Datum}
\end{itemize}
\end{dependencies}

% =========================================================
% Binary64
% =========================================================

\begin{definitionbox}{Definition - IEEE 754 Binary64 Format}
\begin{definition}[IEEE 754 Binary64 Format]
\label{def:numerical-analysis-ieee-754-binary64-format}
The \emph{IEEE 754 binary64 format} is the 64-bit binary floating-point interchange format
with the following field layout:

\begin{itemize}
    \item $1$ sign bit,
    \item $11$ exponent bits,
    \item $52$ fraction bits,
    \item exponent bias $1023$.
\end{itemize}

For a normalized finite binary64 datum, the represented value is
\[
(-1)^s (1.f)_2 2^{e-1023}.
\]
\end{definition}
\end{definitionbox}

\begin{remark*}[Standard quantified statement]
\[
\mathsf{Binary64Format}
=\mathsf{IEEEBinaryFormat}(64,1,11,52,1023).
\]
\end{remark*}

\begin{remark*}[Interpretation]
Binary64 is commonly called \emph{double precision}. It is the standard high-precision
floating-point format for much of scientific and engineering computation.
\end{remark*}

\begin{remark*}[Exposition]
\IEEEFloatLayout{0.14}{1}{11}{52}{63}{52}{51}{binary64}
\end{remark*}

\begin{dependencies}
\begin{itemize}
  \item \hyperref[def:numerical-analysis-ieee-754-standard]{IEEE 754 Standard}
  \item \hyperref[def:numerical-analysis-ieee-754-binary-floating-point-datum]{IEEE 754 Binary Floating-Point Datum}
\end{itemize}
\end{dependencies}

% =========================================================
% Binary128
% =========================================================

\begin{definitionbox}{Definition - IEEE 754 Binary128 Format}
\begin{definition}[IEEE 754 Binary128 Format]
\label{def:numerical-analysis-ieee-754-binary128-format}
The \emph{IEEE 754 binary128 format} is the 128-bit binary floating-point interchange format
with the following field layout:

\begin{itemize}
    \item $1$ sign bit,
    \item $15$ exponent bits,
    \item $112$ fraction bits,
    \item exponent bias $16383$.
\end{itemize}

For a normalized finite binary128 datum, the represented value is
\[
(-1)^s (1.f)_2 2^{e-16383}.
\]
\end{definition}
\end{definitionbox}

\begin{remark*}[Standard quantified statement]
\[
\mathsf{Binary128Format}
=\mathsf{IEEEBinaryFormat}(128,1,15,112,16383).
\]
\end{remark*}

\begin{remark*}[Interpretation]
Binary128 is commonly called \emph{quadruple precision}. It provides substantially greater
range and precision than binary64 at substantially greater storage and computational cost.
\end{remark*}

\begin{remark*}[Exposition]
\IEEEFloatLayout{0.08}{1}{15}{112}{127}{112}{111}{binary128}
\end{remark*}

% =========================================================
% Summary
% =========================================================

\begin{remark*}[Exposition]
The four standard IEEE 754 basic binary interchange formats considered here are:
\[
\text{binary16 }(1,5,10),\qquad
\text{binary32 }(1,8,23),\qquad
\text{binary64 }(1,11,52),\qquad
\text{binary128 }(1,15,112),
\]
where each triple records
\[
(\text{sign bits},\text{ exponent bits},\text{ fraction bits}).
\]
\end{remark*}

\begin{dependencies}
\begin{itemize}
  \item \hyperref[def:numerical-analysis-ieee-754-standard]{IEEE 754 Standard}
  \item \hyperref[def:numerical-analysis-ieee-754-binary-floating-point-datum]{IEEE 754 Binary Floating-Point Datum}
\end{itemize}
\end{dependencies}
